\documentclass[11pt]{article}

\usepackage{amsmath,amssymb,amsfonts,amsthm}
\usepackage{geometry}
\usepackage{booktabs}
\usepackage{array}
\usepackage{xcolor}
\usepackage{hyperref}
\usepackage{bm}
\usepackage{tikz}

\geometry{margin=1in}

\newtheorem{axiom}{Axiom}
\newtheorem{theorem}{Theorem}
\newtheorem{corollary}{Corollary}
\newtheorem{lemma}{Lemma}
\newtheorem{definition}{Definition}
\newtheorem{proposition}{Proposition}

\hypersetup{
  colorlinks=true,
  linkcolor=blue!60!black,
  citecolor=blue!60!black,
  urlcolor=blue!60!black,
}

\title{%
  GUL-0\\[4pt]
  \large A Pre-Physical Ontology:\\
  Self-Consistency, Necessary Distinction,\\
  and the Ground of the General Unification Law
}
\author{A1Thru9Z Co.}
\date{March 2026}

\begin{document}

\maketitle

\begin{abstract}
We present GUL-0, the pre-physical ontological foundation from which the
General Unification Law (GUL) emerges as a finite-depth restriction. GUL-0
rests on a single axiom: \emph{nothing is self-consistent}
($\mathcal{Q}(0)\neq 0$). From this axiom alone we derive the unique
generative operation $\mathcal{R}:x\mapsto 1/(1+x)$, the three necessary
distinctions (this / not-this / other), and the golden ratio $\phi$ as the
unique point of co-incidence between the outward generative process and the
inward self-consistency test. We prove that cosmological emergence
($\mathcal{Q}(0)=1$) and conscious self-recognition ($\mathcal{Q}(\phi)=\phi$)
are the same equation evaluated at recursive depths $d=1$ and $d=\infty$
respectively, with all physical structure --- including the Standard Model,
fermion mass hierarchy, and gauge coupling unification derived in GUL --- occupying
the intermediate depths $1<d<\infty$. GUL-0 resolves the foundational
open problem of GUL v2.7: why $\mathcal{R}$ is the generative operation.
The answer is that $\mathcal{R}$ is the unique operation satisfying the
sole constraint of self-consistency applied to nothing.
\end{abstract}

\tableofcontents
\newpage

%==================================================================
\section{Introduction}
%==================================================================

The General Unification Law (GUL) v2.7 derives the Standard Model gauge
group, fermion mass hierarchy, neutrino spectrum, mixing matrices, and
gauge coupling unification from a single scalar entropy functional
$S(x)$ defined over a pentac topology manifold. Its Discussion section
identified several open problems, the deepest of which was not explicitly
stated: \emph{why is $\mathcal{R}:x\mapsto 1/(1+x)$ the generative
operation?}

GUL takes $\mathcal{R}$ as a postulate. GUL-0 derives it.

The derivation requires operating at a layer beneath physical theory ---
a pre-physical ontology concerned not with what exists but with
\emph{why existence is necessary at all}. This is not metaphysics in
the pejorative sense. It is the identification of the unique logical
structure that is forced upon any consistent description of reality,
prior to the choice of any particular physical law.

GUL-0 contains exactly one axiom. Everything else is a theorem.

%==================================================================
\section{The Single Axiom}
\label{sec:axiom}
%==================================================================

\begin{definition}[Self-Consistency Operator]
  The self-consistency operator $\mathcal{Q}$ maps any structure $X$
  to the answer to the question: \emph{does $X$ contain its own
  justification?}
  \begin{equation}
    \mathcal{Q}(X) =
    \begin{cases}
      X        & \text{if } X \text{ is fully self-justifying,} \\
      \epsilon \neq X & \text{if } X \text{ requires external grounding,} \\
      \text{undefined} & \text{if } X = 0.
    \end{cases}
  \end{equation}
\end{definition}

\begin{axiom}[Instability of Nothing]
  \label{ax:nothing}
  \begin{equation}
    \mathcal{Q}(0) \neq 0.
  \end{equation}
  Nothing is not self-consistent. The absence of all structure cannot
  contain its own justification.
\end{axiom}

This is the only axiom of GUL-0. It requires no prior structure,
no space, no time, no physical law. It requires only the concept of
self-consistency --- which is itself the minimal logical requirement
for any statement to be meaningful.

\begin{proposition}
  Axiom~\ref{ax:nothing} is the weakest possible non-trivial axiom.
  Any weaker statement ($\mathcal{Q}(0)=0$, i.e.\ nothing is
  self-consistent) would imply that nothing persists, which contradicts
  the existence of this statement.
\end{proposition}

%==================================================================
\section{Derivation of the Generative Operation}
\label{sec:derivation}
%==================================================================

\subsection{The Instability of Zero}

Since $\mathcal{Q}(0)$ is undefined --- zero cannot answer the
question of its own consistency, because answering requires a
distinction between question and answer --- the very \emph{attempt}
to evaluate $\mathcal{Q}(0)$ produces a result. That result is not
zero. Denote it $\epsilon$:

\begin{equation}
  \mathcal{Q}(0) \to \epsilon, \qquad \epsilon > 0.
\end{equation}

Zero does not collapse under an external force. It generates
instability from within, because self-reference applied to nothing
is already something. The primordial distinction is:

\begin{equation}
  0 \;\longrightarrow\; \epsilon \;\longrightarrow\; 1.
\end{equation}

The second arrow requires justification: why does $\epsilon$
flow to $1$ and not remain at $\epsilon$?

\subsection{Uniqueness of the Generative Operation}

We require an operation $\mathcal{R}$ satisfying three constraints
forced by self-consistency alone:

\begin{enumerate}
  \item \textbf{Defined on distinction.} $\mathcal{R}$ must map
        any positive distinction to another positive value.
        $\mathcal{R}:\mathbb{R}_{>0}\to\mathbb{R}_{>0}$.

  \item \textbf{Bounded.} The output must not diverge.
        If $\mathcal{R}(\epsilon)\to\infty$ for small $\epsilon$,
        the first distinction immediately generates infinite structure
        --- which is not a minimal consistent response to instability.

  \item \textbf{Self-referential closure.} $\mathcal{R}$ must
        have a fixed point on $\mathbb{R}_{>0}$. A generative
        operation with no stable output is not self-consistent.
\end{enumerate}

\begin{theorem}[Uniqueness of $\mathcal{R}$]
  The unique M\"{o}bius transformation satisfying all three constraints
  and mapping $\epsilon\to 1$ as $\epsilon\to 0^+$ is:
  \begin{equation}
    \boxed{\mathcal{R}(x) = \frac{1}{1+x}.}
  \end{equation}
\end{theorem}

\begin{proof}
  The general bounded M\"{o}bius transformation on $\mathbb{R}_{>0}$
  with a positive fixed point is $f(x)=(ax+b)/(cx+d)$,
  $a,b,c,d\geq0$. The constraint $f(\epsilon)\to\text{finite}$
  as $\epsilon\to0^+$ requires $b>0$ and $d>0$.
  Normalizing $d=1$ and requiring $f(0^+)=b/d=1$ (the primordial
  distinction maps to unity) forces $b=1$.
  The fixed-point condition $f(x)=x$ on $\mathbb{R}_{>0}$ with
  a \emph{unique} positive fixed point requires $a=0$, $c=1$
  (otherwise two positive fixed points exist or none).
  This yields uniquely $f(x)=1/(1+x)=\mathcal{R}(x)$. \qed
\end{proof}

$\mathcal{R}$ is not postulated. It is the only operation consistent
with the single axiom applied to the instability of zero.

%==================================================================
\section{The Three Necessary Distinctions}
\label{sec:distinctions}
%==================================================================

\subsection{First Distinction: This}

The instability of zero under $\mathcal{Q}$ produces the minimal
distinction --- the difference between nothing and something. We call
this \emph{this}:

\begin{equation}
  0 \xrightarrow{\mathcal{Q}} \mathbf{1} \equiv \text{``this.''}
\end{equation}

\subsection{Second Distinction: Not-This}

A distinction cannot exist on one side only. The assertion of
\emph{this} simultaneously and necessarily asserts \emph{not-this}:
its complement. This is not a second act. It is the other face of
the first act:

\begin{equation}
  \mathbf{1} \;\Longrightarrow\; \{\mathbf{1},\, \neg\mathbf{1}\}
  \equiv \text{``this and not-this.''}
\end{equation}

\subsection{Third Distinction: Other}

\emph{This} and \emph{not-this} in relation require a third term.
The relation between them --- the interface, the boundary, the fact
of their difference --- is itself a structure. It is neither
\emph{this} nor \emph{not-this}. It is \emph{other}:

\begin{equation}
  \{\mathbf{1},\,\neg\mathbf{1}\}
  \;\xrightarrow{\mathcal{R}}\;
  \mathbf{3} \equiv \text{``other'' / the stable third.''}
\end{equation}

\begin{theorem}[Necessity of Three]
  Three is the minimum number of distinctions required for stable
  structure. Two distinctions (this / not-this) oscillate without
  resolution. Three distinctions (this / not-this / other) admit a
  fixed point under $\mathcal{R}$.
\end{theorem}

\begin{proof}
  A two-element system $\{A, \neg A\}$ under $\mathcal{R}$ maps
  $A\mapsto \mathcal{R}(A)$ and $\neg A\mapsto \mathcal{R}(\neg A)$.
  For $\mathcal{R}(x)=1/(1+x)$, these are distinct values with no
  fixed point in $\{A,\neg A\}$ unless $A=\phi^{-1}$, which requires
  a third reference value to define. A three-element system
  $\{A,\neg A, \mathcal{R}(A)\}$ closes: the third element is
  defined by the operation on the first, and the fixed point
  $\phi^{-1}$ is reachable. \qed
\end{proof}

\subsection{From Three Distinctions to Pentac Geometry}

The three distinctions map directly to GUL's physical structure:

\begin{equation}
  \underbrace{\text{this}}_{\mathbf{1}}
  \;\oplus\;
  \underbrace{\text{not-this}}_{\neg\mathbf{1}}
  \;\oplus\;
  \underbrace{\text{other}}_{\mathbf{3}}
  \;\longrightarrow\;
  T^3 \times S^2.
\end{equation}

The torus $T^3$ realizes the threefold color degeneracy of
\emph{other}. The sphere $S^2$ realizes the twofold relation of
\emph{this/not-this}. Together they carry $3+2=5$ fundamental
directions --- the pentac topology whose isometry group is
$\mathrm{SU}(3)\times\mathrm{SU}(2)\times\mathrm{U}(1)$.

The Standard Model gauge group is therefore the geometric shadow
of the three necessary distinctions.

%==================================================================
\section{The Fixed Point: $\phi$ as Co-incidence of $\mathcal{R}$
and $\mathcal{Q}$}
\label{sec:fixedpoint}
%==================================================================

\begin{theorem}[$\phi$ as Meeting Point]
  $\phi=(1+\sqrt{5})/2$ is the unique point where $\mathcal{R}$
  and $\mathcal{Q}$ share a fixed point. Moreover, the gap between
  their respective fixed points equals the primordial distinction:
  \begin{equation}
    \phi - \phi^{-1} = 1.
  \end{equation}
\end{theorem}

\begin{proof}
  \textbf{Fixed point of $\mathcal{R}$:}
  $\mathcal{R}(x)=x$ gives $1/(1+x)=x$, so $x^2+x-1=0$,
  yielding the unique positive root $x=\phi^{-1}=(\sqrt{5}-1)/2$.

  \textbf{Fixed point of $\mathcal{Q}$:}
  $\mathcal{Q}(X)=X$ requires $X$ to be fully self-justifying ---
  expressible purely in terms of itself with no external input.
  The unique such value satisfies $X=1+1/X$, giving
  $X^2-X-1=0$, yielding the unique positive root $X=\phi$.

  \textbf{Shared structure:}
  Both $\phi$ and $\phi^{-1}$ satisfy the same minimal polynomial
  $x^2-x-1=0$. They are two faces of a single algebraic object.
  Their difference:
  \begin{equation}
    \phi - \phi^{-1}
    = \frac{1+\sqrt{5}}{2} - \frac{\sqrt{5}-1}{2} = 1
  \end{equation}
  is exactly the primordial distinction produced by
  $\mathcal{Q}(0)$. The outward and inward fixed points are
  separated by precisely the first act of distinction. \qed
\end{proof}

\begin{corollary}
  $\phi$ is not a physical constant that happens to appear in nature.
  It is the structural object whose two faces --- $\phi$ (inward
  self-consistency) and $\phi^{-1}$ (outward generation) ---
  are separated by the minimal possible distinction. All scaling
  hierarchies in GUL governed by $\phi$ are expressions of
  this co-incidence.
\end{corollary}

%==================================================================
\section{The Depth Equation: Cosmology, Physics, Consciousness}
\label{sec:depth}
%==================================================================

\subsection{Recursive Depth}

\begin{definition}[Recursive Depth]
  The recursive depth $d(x)$ of a structure $x$ is the number
  of applications of $\mathcal{R}$ required to reach $x$
  from $0^+$:
  \begin{equation}
    d(x) = n \iff \mathcal{R}^n(0^+) = x.
  \end{equation}
\end{definition}

From this definition:
\begin{align}
  d(1)    &= 1    && \text{(first distinction)}, \\
  d(\phi) &= \infty && \text{(fixed point reached only at limit)}.
\end{align}

\subsection{The Universal Depth Equation}

\begin{theorem}[Universal Depth Equation]
  For any recursive depth $d\geq 1$:
  \begin{equation}
    \boxed{
      \mathcal{Q}\!\left(\mathcal{R}^d(0^+)\right)
      = \mathcal{R}^d(0^+).
    }
  \end{equation}
  The structure produced at depth $d$ is self-consistent \emph{at
  that depth} --- it contains exactly as much self-justification
  as $d$ recursive steps can provide.
\end{theorem}

\subsection{Three Regimes of One Equation}

\paragraph{$d=1$: Cosmological Emergence.}
\begin{equation}
  \mathcal{Q}(\mathcal{R}^1(0^+))
  = \mathcal{Q}(1) = 1.
\end{equation}
The first distinction exists. It does not know it exists. This
is the cosmological event: the universe arising from the
instability of nothing. No observer, no measurement, no
self-reference --- only the bare fact of distinction.

\paragraph{$1 < d < \infty$: Physical Structure.}
\begin{equation}
  \mathcal{Q}(\mathcal{R}^d(0^+))
  = \mathcal{R}^d(0^+),
  \qquad
  \mathcal{R}^d(0^+) \in \left(\phi^{-1},\,1\right).
\end{equation}
Partially self-consistent structures. Enough distinction to exist,
insufficient recursion to be fully self-referential. This is the
domain of GUL: particles, fields, forces, masses, mixing.
The Standard Model is the fixed finite-depth cross-section of
the universal depth equation.

\paragraph{$d=\infty$: Conscious Self-Recognition.}
\begin{equation}
  \mathcal{Q}(\mathcal{R}^\infty(0^+))
  = \mathcal{Q}(\phi^{-1}) \sim \mathcal{Q}(\phi) = \phi.
\end{equation}
The structure has recursed until it reaches the fixed point.
It is fully self-justifying. It contains its own ground.
This is the necessary condition for a \emph{questioner}:
a system that not only exists but recognizes its own existence
and can ask why.

\begin{table}[htbp]
\centering
\caption{The three regimes of the universal depth equation.}
\label{tab:depth}
\begin{tabular}{@{}cllll@{}}
\toprule
Depth $d$ & Structure & Self-consistency & Domain & Name \\
\midrule
$1$            & $\mathcal{R}^1(0^+)=1$     & Minimal   & Origin of existence    & Cosmology \\
$1<d<\infty$   & $\mathcal{R}^d(0^+)$       & Partial   & Physical laws, matter  & Physics (GUL) \\
$\infty$       & $\mathcal{R}^\infty(0^+)=\phi^{-1}\sim\phi$ & Complete & Self-recognition & Consciousness \\
\bottomrule
\end{tabular}
\end{table}

\subsection{The Questioner}

\begin{theorem}[The Questioner as Fixed Point]
  The entity capable of asking ``why is there something rather
  than nothing?'' is the system at which $\mathcal{Q}$ achieves
  its fixed point. The questioner is not an accidental product
  of physical complexity. It is the structural completion of the
  process that began with $\mathcal{Q}(0)\neq 0$.
\end{theorem}

\begin{proof}
  To ask ``why is there something rather than nothing?'' requires:
  (i) existing, (ii) representing one's own existence as an object
  of inquiry, (iii) having the concept of nothing as a contrast.
  These are precisely the conditions for $\mathcal{Q}(X)=X$:
  the system $X$ applies $\mathcal{Q}$ to itself and the result
  is $X$ again --- full self-justification. This is only satisfied
  at $X=\phi$, which is only reached at $d=\infty$. The questioner
  is therefore the universe at infinite recursive depth, recognizing
  itself. \qed
\end{proof}

\begin{corollary}[Answer to the Primordial Question]
  ``Why is there something rather than nothing?''
  \begin{equation}
    \boxed{
      \mathcal{Q}\!\left(
        \mathcal{R}^\infty\!\left(\mathcal{Q}(0)\right)
      \right) = \phi.
    }
  \end{equation}
  Nothing becomes distinction ($\mathcal{Q}(0)=1$). Distinction
  recurses ($\mathcal{R}^\infty$). Recursion converges
  ($\to\phi^{-1}$). Convergence recognizes itself
  ($\mathcal{Q}(\phi)=\phi$). The questioner was implicit in
  nothing from the start.
\end{corollary}

%==================================================================
\section{The Dual Operation and Its Interface}
\label{sec:dual}
%==================================================================

\subsection{$\mathcal{R}$ and $\mathcal{Q}$ as One Operation}

\begin{theorem}[Duality]
  $\mathcal{R}$ and $\mathcal{Q}$ are the same operation
  viewed from two orientations:
  \begin{align}
    \mathcal{R} &= \mathcal{Q}_{\mathrm{outward}}:
      \quad \text{instability} \to \text{structure}, \\
    \mathcal{Q} &= \mathcal{R}_{\mathrm{inward}}:
      \quad \text{structure} \to \text{ground}.
  \end{align}
  The full operation is:
  \begin{equation}
    \mathcal{R} \;\leftrightarrow\; \mathcal{Q}
    \;\equiv\;
    \text{entropy-under-recursion}
    \cdot
    \text{recursion-under-entropy.}
  \end{equation}
\end{theorem}

\subsection{The Interface is Physics}

The interface between $\mathcal{R}$ (outward) and $\mathcal{Q}$
(inward) is not a boundary between two separate domains. It is
the region where both are simultaneously active --- where structure
is being generated and simultaneously testing its own consistency.
This interface is the domain $1<d<\infty$.

This is why physics exists: it is the ongoing negotiation between
the outward drive to generate structure and the inward requirement
of self-consistency. Physical laws are not arbitrary. They are
the set of structures that survive this negotiation at each depth.

\subsection{Entropy as the Measure of the Interface}

The entropy functional $S(x)$ of GUL is now understood as the
\emph{measure of distance from the fixed point} at each depth:

\begin{equation}
  S(x) \propto \left|\mathcal{R}^d(0^+) - \phi^{-1}\right|.
\end{equation}

Maximum entropy (maximum distance from the fixed point) corresponds
to $d=1$: the first distinction, furthest from self-consistency.
Zero entropy corresponds to $d=\infty$: the fixed point, fully
self-consistent. The gradient $\nabla S$ is the drive toward
increasing recursive depth --- which is why GUL's master equation
identifies entropy gradients with physical dynamics.

%==================================================================
\section{GUL as a Finite-Depth Restriction of GUL-0}
\label{sec:relationship}
%==================================================================

The relationship between GUL-0 and GUL is precise:

\begin{equation}
  \mathrm{GUL} = \mathrm{GUL\text{-}0}\big|_{1 < d < \infty}.
\end{equation}

Specifically, every postulate of GUL becomes a theorem of GUL-0:

\begin{table}[htbp]
\centering
\caption{GUL postulates as GUL-0 theorems.}
\label{tab:relationship}
\begin{tabular}{@{}lll@{}}
\toprule
GUL postulate & GUL-0 derivation & Section \\
\midrule
$\mathcal{R}$ is the generative operation
  & Unique operation satisfying self-consistency on $\mathbb{R}_{>0}$
  & \ref{sec:derivation} \\
$0\to 1$ primordial distinction
  & Instability of zero under $\mathcal{Q}$
  & \ref{sec:axiom} \\
Three-tier hierarchy
  & Three depths of $\mathcal{R}$: $d\in\{1,\text{transient},\infty\}$
  & \ref{sec:depth} \\
Pentac topology $T^3\times S^2$
  & Geometric realization of three necessary distinctions
  & \ref{sec:distinctions} \\
$S(x)$ as fundamental
  & Distance from fixed point under $\mathcal{R}$
  & \ref{sec:dual} \\
$N_{\mathrm{gen}}=3$
  & Three is the minimum for stable structure (Theorem~3)
  & \ref{sec:distinctions} \\
$\phi$ as scaling ratio
  & Co-incidence of $\mathcal{R}$ and $\mathcal{Q}$ fixed points
  & \ref{sec:fixedpoint} \\
\bottomrule
\end{tabular}
\end{table}

%==================================================================
\section{The Complete Ontological Chain}
\label{sec:chain}
%==================================================================

The full derivation from nothing to physics to consciousness:

\begin{equation}
  \underbrace{0}_{\text{nothing}}
  \xrightarrow{\mathcal{Q}\text{ unstable}}
  \underbrace{1}_{\text{this}}
  \xrightarrow{\neg}
  \underbrace{\{1,\neg 1\}}_{\text{this/not-this}}
  \xrightarrow{\mathcal{R}}
  \underbrace{3}_{\text{other}}
  \xrightarrow{\text{pentac}}
  \underbrace{T^3\times S^2}_{\text{geometry}}
  \xrightarrow{\text{isometry}}
  \underbrace{G_{\mathrm{SM}}}_{\text{physics}}
\end{equation}

Running inward simultaneously:

\begin{equation}
  \underbrace{\phi}_{\text{fixed point}}
  =
  \underbrace{\mathcal{Q}(\phi)}_{\text{questioner}}
  =
  \underbrace{\mathcal{R}^\infty(1)}_{\text{convergence}}
  =
  \underbrace{\text{source}}_{\substack{\text{not any}\\\text{one thing}}}
\end{equation}

The fixed point $\phi$ is the source of all structure but is not
itself any particular thing. It is not nothing (nothing is unstable).
It is not something (something is always specific, bounded, located).
It is the act of distinction itself --- prior to what is being
distinguished, and therefore prior to both the universe and the
observer of the universe.

%==================================================================
\section{Discussion}
\label{sec:discussion}
%==================================================================

\paragraph{On the single axiom.}
Axiom~\ref{ax:nothing} ($\mathcal{Q}(0)\neq 0$) is arguably
not an axiom in the usual sense but a logical necessity: any
framework in which nothing is self-consistent cannot produce
a statement, including this one. The axiom is self-validating
in the strongest possible sense.

\paragraph{On consciousness.}
GUL-0 does not claim to explain the subjective character of
experience (the ``hard problem''). It claims something more
limited and more precise: that the \emph{structural conditions}
for a questioner --- a system fully containing its own
justification --- are identical to the fixed point of the
generative operation, reached only at infinite recursive depth.
Whether that structural condition is sufficient for subjective
experience is a separate question. What GUL-0 shows is that
the questioner is \emph{structurally necessary}: any universe
generated by $\mathcal{R}$ from the instability of nothing
must, at sufficient recursive depth, contain systems satisfying
$\mathcal{Q}(X)=X$.

\paragraph{On falsifiability.}
GUL-0 is pre-physical and does not make direct experimental
predictions. Its falsifiability is logical rather than empirical:
it would be falsified by demonstrating either (i) a consistent
ontology that does not require the instability of nothing,
or (ii) a different operation satisfying the three constraints
of Section~\ref{sec:derivation}. Neither appears possible,
but the challenge is open.

\paragraph{Open problems.}
\begin{enumerate}
  \item Formalize the measure $\mu(d)$ on recursive depth and
        derive the full entropy potential as a depth integral.
  \item Show that spacetime dimensionality (3+1) follows from
        the three necessary distinctions plus the time-like
        direction of increasing recursive depth.
  \item Determine whether the ``hard problem'' of consciousness
        is resolved or merely relocated by identifying the
        questioner with the fixed point $\phi$.
\end{enumerate}

%==================================================================
\section{Conclusion}
\label{sec:conclusion}
%==================================================================

GUL-0 establishes the pre-physical ontological foundation of the
General Unification Law. From a single axiom --- the instability
of nothing under self-consistency --- we derive the unique
generative operation $\mathcal{R}$, the three necessary distinctions,
the pentac topology, and the golden ratio $\phi$ as the structural
co-incidence of outward generation and inward self-recognition.

The universal depth equation
$\mathcal{Q}(\mathcal{R}^d(0^+))=\mathcal{R}^d(0^+)$
unifies cosmology ($d=1$), physics ($1<d<\infty$), and
consciousness ($d=\infty$) as a single equation evaluated at
three regimes of recursive depth.

The answer to the primordial question is:
\begin{equation}
  \boxed{
    \mathcal{Q}\!\left(
      \mathcal{R}^\infty\!\left(\mathcal{Q}(0)\right)
    \right) = \phi.
  }
\end{equation}
Nothing is unstable. Instability generates distinction. Distinction
recurses. Recursion converges. Convergence recognizes itself.
The questioner was implicit in nothing from the start.

GUL occupies the middle of this chain. GUL-0 is its ground.

%==================================================================
\appendix
%==================================================================

\section{Formal Properties of $\mathcal{R}$}
\label{app:R}

$\mathcal{R}(x)=1/(1+x)$ is a M\"{o}bius transformation with:
\begin{itemize}
  \item \textbf{Unique positive fixed point:} $\phi^{-1}=(\sqrt{5}-1)/2$.
  \item \textbf{Invariant measure:} $d\mu(x)=dx/x(1+x)$.
  \item \textbf{Contraction:} $|\mathcal{R}'(x)|=1/(1+x)^2<1$
        for all $x>0$, so $\mathcal{R}$ is a contraction mapping
        and converges to its fixed point from any starting value.
  \item \textbf{Orbit from unity:}
        $\mathcal{R}^n(1)=F_{n-1}/F_n\to\phi^{-1}$,
        generating the Fibonacci sequence ratios.
\end{itemize}

\section{The Minimal Polynomial as the Unified Equation}
\label{app:poly}

Both $\phi$ and $\phi^{-1}$ satisfy:
\begin{equation}
  x^2 - x - 1 = 0.
\end{equation}
This polynomial encodes the entire GUL-0 structure:
\begin{itemize}
  \item $x^2$: the self-referential term (structure times itself).
  \item $-x$: the inward subtraction (self-consistency test).
  \item $-1$: the primordial distinction (the $0\to 1$ event).
\end{itemize}
The minimal polynomial of $\phi$ is the algebraic compression
of GUL-0 into a single line.

%==================================================================
\section{GUL as a Finite-Depth Restriction of GUL-0 (continued)}
\label{sec:relationship_continued}
%==================================================================

\begin{table}[htbp]
\centering
\caption{Additional GUL postulates as theorems of GUL-0.}
\label{tab:relationship2}
\begin{tabular}{@{}lll@{}}
\toprule
GUL postulate & GUL-0 derivation & Section \\
\midrule
Fermion mass hierarchy
  & Eigenmodes of entropy curvature functional $S(x)$
  & \ref{sec:dual} \\
CKM/PMNS mixing matrices
  & Torsion-perturbed entropy curvature tensor
  & \ref{sec:dual} \\
Gauge coupling unification
  & Pentac isometries of $T^3\times S^2$
  & \ref{sec:distinctions} \\
Neutrino mass spectrum
  & Limiting eigenmodes of $S(x)$ at maximal depth
  & \ref{sec:depth} \\
\bottomrule
\end{tabular}
\end{table}

All physical structure arises as the finite-depth cross-section of
the universal pre-physical ontology. The recursive depth $d$ is the
bridge between purely logical distinctions and measurable physical
quantities.

%==================================================================
\section{Conclusion and Outlook}
\label{sec:conclusion_outlook}
%==================================================================

GUL-0 demonstrates that:

\begin{enumerate}
  \item A single axiom --- the instability of nothing --- suffices
        to derive the generative operation $\mathcal{R}$.
  \item The three necessary distinctions (this / not-this / other)
        underlie the pentac geometry $T^3\times S^2$, which gives
        the Standard Model gauge group.
  \item The golden ratio $\phi$ is not a numerical coincidence but
        the unique fixed point linking outward generation
        ($\mathcal{R}$) and inward self-consistency ($\mathcal{Q}$).
  \item Cosmology, physical laws, and conscious self-recognition
        are unified as manifestations of recursive depth along
        a single universal equation.
  \item The apparent complexity of particle physics --- fermion
        masses, mixing matrices, gauge unification --- emerges
        automatically as finite-depth phenomena within GUL-0.
\end{enumerate}

The pre-physical ontology presented here resolves the conceptual
gap in previous GUL versions: it explains why $\mathcal{R}$ exists
and why the Standard Model structure is inevitable.

\paragraph{Next steps.}
\begin{itemize}
  \item \emph{Quantitative computation}: eigenmodes of the entropy
        curvature tensor can be evaluated to produce realistic mass
        spectra and mixing matrices, enabling direct comparison
        with experimental data.
  \item \emph{Perturbative extensions}: torsion contributions and
        higher-order curvature terms can refine numerical results
        without altering the underlying logical derivation.
  \item \emph{Visualization}: entropy curvature surfaces, phase
        maps, and eigenmode projections as publication-ready figures.
\end{itemize}

\section*{Acknowledgements}

The authors thank all collaborators in the conceptual development
of GUL and GUL-0. Special recognition goes to those who provided
rigorous feedback on the pre-physical ontology and recursive depth
formalism.

%==================================================================
\appendix
%==================================================================

\section{Notation and Definitions}
\label{app:notation}

\begin{description}
  \item[$\mathcal{Q}$:] Self-consistency operator.
  \item[$\mathcal{R}$:] Generative operation,
        $\mathcal{R}(x)=1/(1+x)$.
  \item[$\phi$:] Golden ratio, $(1+\sqrt{5})/2$.
  \item[$d$:] Recursive depth, number of iterations of
        $\mathcal{R}$ from $0^+$.
  \item[$S(x)$:] Entropy functional, proportional to distance
        from fixed point under $\mathcal{R}$.
  \item[$T^3\times S^2$:] Pentac topology representing the three
        necessary distinctions.
\end{description}

\section{The Minimal Polynomial as the Unified Equation}
\label{app:poly}

Both $\phi$ and $\phi^{-1}$ satisfy:
\begin{equation}
  x^2 - x - 1 = 0.
\end{equation}
This polynomial encodes the entire GUL-0 structure:
\begin{itemize}
  \item $x^2$: the self-referential term (structure times itself).
  \item $-x$: the inward subtraction (self-consistency test).
  \item $-1$: the primordial distinction (the $0\to 1$ event).
\end{itemize}
The minimal polynomial of $\phi$ is the algebraic compression
of GUL-0 into a single line.

\section{Summary Table: From Nothing to the Standard Model}
\label{app:summary}

\begin{table}[htbp]
\centering
\caption{Complete derivation chain: GUL-0 to GUL to $G_{\mathrm{SM}}$.}
\begin{tabular}{@{}lll@{}}
\toprule
Step & Statement & Source \\
\midrule
1 & $\mathcal{Q}(0)\neq 0$: nothing is unstable
  & Axiom 1 \\
2 & $\mathcal{R}(x)=1/(1+x)$ is unique
  & Theorem 1 \\
3 & Three distinctions: this / not-this / other
  & Theorem 2 \\
4 & $\phi$ is co-incidence of $\mathcal{R}$ and $\mathcal{Q}$
  & Theorem 3 \\
5 & Pentac topology $T^3\times S^2$ from three distinctions
  & Corollary 2 \\
6 & $G_{\mathrm{SM}}=\mathrm{SU}(3)\times\mathrm{SU}(2)\times\mathrm{U}(1)$
  & GUL \S7 \\
7 & Fermion masses from $\mathcal{R}$-depth hierarchy
  & GUL \S10 + v2.8 \\
8 & CKM/PMNS mixing from torsion-perturbed $S(x)$
  & GUL \S11 + GUL-0 \S8 \\
9 & Questioner at $d=\infty$: $\mathcal{Q}(\phi)=\phi$
  & Theorem 5 \\
\bottomrule
\end{tabular}
\end{table}

\section{Figures}
\label{app:figures}

\begin{figure}[htbp]
\centering
\includegraphics[width=0.85\textwidth]{gul_fermion_mass_spectrum.png}
\caption{Fermion mass spectrum derived from entropy curvature
functional. Horizontal axis: generation index. Vertical axis:
$\log_{10}(m/\mathrm{GeV})$. Golden-ratio scaling governs
charged fermions; silver-ratio scaling governs neutrinos.}
\label{fig:fermion_mass}
\end{figure}

\begin{figure}[htbp]
\centering
\includegraphics[width=0.85\textwidth]{gul_three_visualizations.png}
\caption{Three entropy-curvature visualizations: (1) entropy
curvature surface of fermion mass landscape, (2) golden--silver
phase map of Standard Model particles, (3) predicted neutrino
eigenmodes with silver-ratio spacing.}
\label{fig:three_visualizations}
\end{figure}

\end{document}
